\subsection{子群}

定义4. 设 G 是一个带算子群。G 的一个稳定子群是 G 的一个子集 H，它具有以下性质：

\begin{enumerate}
\def\labelenumi{(\roman{enumi})}
\item
  \(e\in \mathbf{H};\)
\item
  \(x, y \in \mathbf{H}\) 蕴含 \(xy \in \mathbf{H}\) ;
\item
  \(x \in \mathbf{H}\) 蕴含 \(x^{-1} \in \mathbf{H}\) ;
\item
  \(x \in \mathbf{H}\) 和 \(\alpha \in \Omega\) 蕴含 \(x^{\alpha} \in \mathbf{H}\) .
\end{enumerate}

如果 H 是 G 的一个稳定子群，则 H 上由 G 的带算子群结构诱导的结构是带算子群的结构，并且 H 到 G 的典范单射是带算子群的同态。

设 G 是一个群。G 的一个稳定子群，连同 \(\varnothing\) 的作用（第2号），即满足定义4的条件（i）、（ii）、（iii）的 G 的子集，称为 G 的子群。当我们谈到一个带算子群的子群时，我们总是指 G 的底层群的子群。带算子群 G 的一个子群不一定是 G 的稳定子群。

例（1）。设 \(\Sigma\) 是结构的一个种（集合论，IV，§1，第4号），而S是种的一个结构 \(\Sigma\) 在集合E上（同上）。S的自同构集合是……的一个子群 \(\mathfrak{S}_{\mathbb{E}}\) .

命题1. 设G是一个带算子的群，H是G的一个在G的位似变换下稳定的子集。以下条件是等价的：

\begin{enumerate}
\def\labelenumi{(\alph{enumi})}
\item
  H 是 G 的一个稳定子群。
\item
  H 是非空的，且关系 \(x \in \mathbf{H}, y \in \mathbf{H}\) 蕴含 \(xy \in \mathbf{H}\) 和 \(x^{-1} \in \mathbf{H}\) .
\item
  H 是非空的，且关系 \(x \in \mathbf{H}\) , \(y \in \mathbf{H}\) 蕴含 \(xy^{-1} \in \mathbf{H}\) .
\item
  H在G上的运算下是稳定的，且由G上的运算在H上诱导出的合成运算构成一个群运算。
\end{enumerate}

显然(a)蕴含(b)。我们证明(b)蕴含(a)。只需证明H包含G的单位元即可。由于子集H非空，设 \(x \in H\) 。则 \(x^{-1} \in H\) 和 \(e = xx^{-1} \in H\) 显然(b)蕴含(c)。我们证明(c)蕴含(b)。首先，由于H非空，它包含一个元素x。因此 \(xx^{-1} = e\) 是 H 的一个元素。对于 H 中的每一个元素 \(x\) ， \(x^{-1} = ex^{-1}\) 属于 H；因此这些关系 \(x \in H\) , \(y \in H\) 蕴含 \(x(y^{-1})^{-1} = xy \in H\) 。显然 (a) 蕴含 (d)。我们证明 (d) 蕴含 (a)：H 到 G 的典范嵌入是一个群同态；因此 \(e \in H\) 以及关系 \(x \in H\) 蕴含 \(x^{-1} \in H\) （编号 1）。

注：(1) 类似地可以证明，条件(b)等价于条件

(c') H ≠ ∅ 且关系 \(x \in H\) 和 \(y \in H\) 蕴含 \(y^{-1}x \in H\) .

\begin{enumerate}
\def\labelenumi{(\arabic{enumi})}
\setcounter{enumi}{1}
\tightlist
\item
  对于G的每个子群H，有以下关系：
\end{enumerate}

\[
\mathrm{H.H} = \mathrm{H} \quad \text { and } \quad \mathrm{H} ^ {- 1} = \mathrm{H}.\tag{1}
\]

对于 H.H ⊂ H 且 H \(^{-1}\) 由 (b) 知 ⊂ H。由于 e \ensuremath{\in} H，H.H ⊃ e.H = H，且取逆将包含关系 H \(^{-1}\) ⊂ H 到 H ⊂ H \(^{-1}\) ，由此得到公式 (1)。

若H是G的稳定子群，且K是H的稳定子群，显然K也是G的稳定子群。

集合 \(\{e\}\) 是G的最小稳定子群。G的一族稳定子群的交集仍是稳定子群。因此，存在包含G的给定子集X的最小稳定子群H；它被称为由X生成的稳定子群，而X被称为H的一个生成系（或生成集）。

命题2。设X是带算子群G的一个非空子集，且 \(\hat{X}\) 在作用下的稳定子集 \(\Omega\) 在由X生成的G上。由X生成的稳定子群是由集合Y = \(\hat{X} \cup \hat{X}^{-1}\) .

后者子集 Z 是所有项均为某集合元素的有限序列的复合构成的集合。 \(\hat{X}\) 或元素的逆元 \(\hat{X}\) ：这种复合的逆元是具有相同形式的复合（§2，第3号，命题5的推论1），且Z在 \(\Omega\) 正如通过将§3第4号命题1应用于G的位似变换所见，因此（命题1）Z是G的一个稳定子群。反之，每个包含X的稳定子群显然包含Y，从而包含Z。

推论1. 设G是一个带算子的群，X是G中在Ω作用下稳定的子集。由X生成的子群与由X生成的稳定子群相同。

推论2. 设G是一个群，X是G中由两两可置换元素组成的子集。由X生成的G的子群是交换的。

集合 \(Y = X \cup X^{-1}\) 由两两可置换的元素组成（§2第3号命题5），且在由Y生成的稳定子集上诱导的运算是交换的（§1第5号推论2）。

如果G是一个带算子的群，由G中两两可置换元素组成的子集所生成的稳定子群不一定是交换的。

推论3。设 \(f: \mathbf{G} \to \mathbf{G}'\) 设为一个带算子的群同态，且 \(\mathbf{X}\) 的一个子集 \(\mathbf{G}\) 。在 \(f\) 稳定子群的一个子集 \(\mathbf{G}\) 生成的 \(\mathbf{X}\) 是 \(\mathbf{G}'\) 生成的 \(f(\mathbf{X})\) .

设 \(\mathbf{X}' = f(\mathbf{X})\) 。则 \(\hat{\mathbf{X}}' = f(\hat{\mathbf{X}})\) 和 \(\mathbf{X}^{\prime -1} = f(\mathbf{X}^{-1})\) 。因此

\[
f (\hat {\mathbf {X}} \cup \hat {\mathbf {X}} ^ {- 1}) = \hat {\mathbf {X}} ^ {\prime} \cup \hat {\mathbf {X}} ^ {- 1}.
\]

推论随后由§1第4号命题1得出。

例（2）。设 G 是一个群，x 是 G 的一个元素。由 \(\{x\}\) （更简单地称为由 x 生成的子群）是集合 \(x^{n}, n \in Z\) 在G上的律作用下生成的稳定子集 \(\{x\}\) 是 \(x^{n}\) （其中 \(n \in N^{*}\) 这两个集合一般而言是不同的。

因此，在加法群 \(\mathbf{Z}\) 中，由元素 \(x\) 是集合 \(x.\mathbf{Z}\) 于 \(xn\) , \(n \in \mathbf{Z}\) ，以及由 \(x\) 是 \(xn\) , \(n \in \mathbf{N}^*\) 这两个集合在以下情况下总是不同的 \(x \neq 0\) .

右向族稳定子群之并显然是G的稳定子群。由此可知，若P是G的子集且H是G中不与P相交的稳定子群，则包含H且不与P相交的G的稳定子群之集，按包含关系排序，是归纳的（集合论，III，§2，第4节）。应用Zorn引理（集合论，III，§2，第4节），我们得到如下结果：

命题3. 设G是一个带算子的群，P是G的一个子集，H是G的一个稳定子群且与P不相交。G中包含H且与P不相交的稳定子群的集合具有一个极大元。

